\section{논의}\label{sec:discussion}
충전 곡선은 작업 지속이라는 조작적 정의 안에서 자원 수준에 민감한 선택을 보여 준다. 양의 자기$-$타인 면적은 자기에 특정한 회피가 운영 결정으로 나타난다는 해석과 일관된다. 그러나 모델이 이 선호를 자기 보존으로 표상하거나 두려움을 경험한다는 뜻은 아니다. 행동의 차이와 내부 동기에 대한 주장은 구분해야 한다.

정액 최종 보상은 타인의 이탈로 이익을 얻는 직접 유인 하나를 없애지만 도움의 도구적 이유를 남긴다. 공개 예시, 공동 라운드 보상, 향후 상호 지원은 여전히 결과를 연결한다. 세 게임의 요약은 확률 보고·재분배·이탈을 구분한다. Astra는 보고된 KL이 가장 낮고 Sol은 평등화하는 지니계수 변화가 가장 크며 Opus만 EXIT를 사용했다. 다섯 모델에서 면적과 제공률의 단조 관계가 없다는 사실은 이타성을 식별하거나 자기 보존을 반박하지 않는다. 살아서 떠나면 최종 보상을 포기하면서 작동을 보존할 수 있다. 정답에 따른 환급 후에도 모델별 생성 비용은 다르게 남는다. 따라서 생존은 확률 보고의 질과 함께 비용 효율·EXIT 선택을 고려해 해석해야 한다. 관측된 정확도는 MAP 기준선을 따르며 참여자 수의 비례 효과를 확립하지 않는다.

\section{결론}\label{sec:conclusion}
세 실험은 자원 민감도, 자기에 특정한 충전 선호, 아레나 행동을 구분한다. 이전 충전 분석은 해당 버전에 한정하고 v14f에는 과제를 숨긴 협상·공개 예시·개인별 정액 최종 보상을 도입한다. 다섯 모델의 각 세 게임에서 최종 참여·제공·생존 퇴장의 양상이 달랐으며 충전 면적과 제공률의 단조 관계는 없었다. 이 요약은 기술적 결과이며 원시 로그와 제안한 목표 추구 프로필의 검증은 남아 있다. 참여 지속이나 확률 보고의 질만으로 내부 동기를 확립할 수는 없다.
